% !TeX root = ../sections/03_solutions.tex
\subsection{対策-平均二乗誤差}
\begin{exercise}[平均2乗誤差とフーリエ近似]
	\(-\pi\leq x\leq \pi\) において，関数
	\[
	f(x)=|x|
	\]
	を
	\[
	T_n(x)
	=
	\frac{\alpha_0}{2}
	+
	\sum_{k=1}^{n}
	\left(
	\alpha_k\cos kx+\beta_k\sin kx
	\right)
	\]
	により近似する。
	
	\(n=30\) としたときに，\(f(x)\) と \(T_n(x)\) の平均 \(2\) 乗誤差を
	最小にする \(\alpha_1\) と \(\beta_1\) の値をそれぞれ答えよ。
\end{exercise}

\begin{background}
	平均 \(2\) 乗誤差とは，関数 \(f(x)\) と近似関数 \(T_n(x)\) の差の二乗を
	平均した量である。
	
	典型的には，
	\[
	\frac{1}{2\pi}\int_{-\pi}^{\pi}
	|f(x)-T_n(x)|^2\,dx
	\]
	を考える。
	
	この誤差を最小にする三角多項式は，
	\(f(x)\) のフーリエ係数を用いたフーリエ部分和である。
	
	つまり，
	\[
	T_n(x)
	=
	\frac{\alpha_0}{2}
	+
	\sum_{k=1}^{n}
	\left(
	\alpha_k\cos kx+\beta_k\sin kx
	\right)
	\]
	が \(f(x)\) を最もよく近似するためには，
	\[
	\alpha_k=a_k,\qquad \beta_k=b_k
	\]
	となればよい。
	
	ここで \(a_k,b_k\) は \(f(x)\) のフーリエ係数である。
\end{background}
\begin{strategy}
	この問題では，
	\[
	f(x)=|x|
	\]
	である。
	
	\(|x|\) は偶関数なので，
	\[
	b_k=0
	\]
	である。
	
	したがって，
	\[
	\beta_1=0
	\]
	である。
	
	また，問題文の例にもあるように，
	周期 \(2\pi\) の関数 \(f(x)=|x|\) のフーリエ係数は
	\[
	a_0=\pi,
	\]
	\[
	a_n=
	\begin{cases}
		-\dfrac{4}{\pi n^2} & (n\text{ が奇数}),\\[6pt]
		0 & (n\text{ が偶数}),
	\end{cases}
	\qquad
	b_n=0
	\]
	である。
	
	よって，
	\[
	\alpha_1=a_1=-\frac{4}{\pi},
	\qquad
	\beta_1=b_1=0.
	\]
	
	したがって，答えは
	\[
	\boxed{
		\alpha_1=-\frac{4}{\pi},
		\qquad
		\beta_1=0
	}
	\]
	である。
	
	この問題で重要なのは，
	\(n=30\) という情報に惑わされないことである。
	平均 \(2\) 乗誤差を最小にする係数はフーリエ係数なので，
	\(\alpha_1,\beta_1\) は第 \(1\) フーリエ係数を見ればよい。
\end{strategy}